\section{Simulation against a known answer}
\label{sec:sim}

Everything so far has been measured on real logs, where the answer is not
known. This section measures the estimator where it is. The full tables are in
Supplement~\ref{app:tables} and the diagnosis of the one world that failed is
in Supplement~\ref{app:noisy}.

\subsection{Six worlds, two estimands}
\label{sec:simworlds}

A simulation that fits a logistic data-generating process with logistic
regression tests an unusually favourable arrangement. This one uses \nWorlds\
worlds --- four of them misspecified for the estimator --- at \nSimReps\
replicates each. The covariate space is finite, so the joint distribution of
cell and outcome is enumerated and the population AUC of any scoring rule is
computed in closed form rather than simulated. The worlds are a logistic
surface, a nonlinear one with an interaction and a threshold the logistic
model cannot see, a drifting one, a sparse register with \nSimSparse\ levels,
an imbalanced outcome at prevalence \simPrevalence, and a noisy register in
which a share of values are wrong.

Two estimands are separated because under misspecification they differ.
$V_{\mathrm{oracle}}$ is the increment between the two Bayes-optimal scoring
rules --- what the register is worth in the world. $V_{\mathrm{limit}}$ is the
increment between the two fitted pipelines' limits --- what the estimator
estimates. A resampling interval can only cover the second, and reporting its
coverage of the first as though it were the estimator's coverage is a category
error. Both are reported in Table~\ref{tab:sim31}. Where a register value is
corrupted the cells the generator writes down are not the cells the estimator
observes, so the truth has to be marginalised over the corruption; the two
experiments that found that and the check that guards it are in
Supplement~\ref{app:noisy}.

\subsection{Coverage, bias and width}
\label{sec:simcoverage}

Coverage estimates carry a Monte Carlo standard error of about
$\simCoverageSE$ percentage points at \nSimReps\ replicates, which is the
resolution every comparison here should be read at;
Figure~\ref{fig:coverage} draws them by construction and world.

\textbf{Where the refit helps.} The nested basic interval covers at least as
well as the fixed-model percentile interval on \nWorldsNestedAtLeastNaive\ of
the \nWorldsSim\ worlds, and its median across them is
\coverageNestedBasicMedian\ against the fixed-model interval's
\coverageNaiveMedian. Both fall a little short of nominal on most worlds, by
more than the Monte Carlo resolution.

\textbf{Where the percentile interval fails, and why.} On the sparse world it
collapses to \coverageSparsePct, and the \emph{nested} percentile interval is
worse there than the fixed-model one. The cause is Section~\ref{sec:boot}: a
bootstrap training resample holds about $63\%$ of the distinct register
levels, so every refit sees a sparser register and the bootstrap distribution
shifts below the estimate. The basic interval pivots that shift out and
restores coverage --- \coverageSparseBasic\ --- and costs nothing where there
was nothing to repair, at a median \coverageNestedBasicMedian\ and a minimum
\coverageNestedBasicMin\ across all \nWorlds\ worlds. The bias-corrected
percentile does \emph{not} repair the sparse world (\coverageSparseBc), which
is why every interval and band in this paper uses the basic construction. The
$m$-out-of-$n$ subsampling interval, the standard alternative in exactly this
regime, does not win either: \coverageSparseMofn\ on the sparse world and a
median \coverageMofnMedian\ across the \nWorldsMofnPriced\ worlds it was
priced on.

\textbf{What no interval repairs.} Against $V_{\mathrm{oracle}}$ rather than
$V_{\mathrm{limit}}$, every construction undercovers on the drifting world ---
\coverageDriftOracle\ at best --- and it should: when the coefficients move
between eras the estimator is not consistent for the oracle quantity, and the
largest gap is \biasOracleMax\ in AUC units. Under misspecification an
interval around the estimator's own limit is not an interval around the
quantity a reader cares about, and no resampling scheme repairs that.

\subsection{Sample size, register cardinality and block length}
\label{sec:simsensitivity}

A construction that works at one sample size, one register cardinality and one
block length has been validated at a point, not in a regime. Sample size and
register cardinality are therefore varied \emph{jointly}, up to $K = \simKMax$
levels --- the cardinality of the case study's own register --- and the block
length is varied around its rule of thumb.

\textbf{The ratio $K/n$ captures most of the dependence, and not all of it.}
Across the plane the basic interval covers a median \coverageGridMedian\ and a
minimum \coverageGridMin, the last at a ratio of \simGridWorstRatio. Below
$K/n = \simGridRatioSafe$ its coverage lies between \coverageGridSafeMin\ and
\coverageGridSafeMax\ --- short of nominal everywhere, by about one Monte
Carlo standard error at \nPlaneGridReps\ replicates --- and it has fallen
below three quarters by \simGridRatioBad. The ratio is the right first-order
summary --- the mechanism of Section~\ref{sec:boot} is why, and
varying $n$ alone would not have exposed the regime the corpus occupies ---
but it is not sufficient, and the plane says so in two places. The two cells
closest to each other in ratio --- a register of \simGridSameRatioKa\ levels
at \simGridSameRatioNa\ training rows, and one of \simGridSameRatioKb\ at
\simGridSameRatioNb\ --- cover \simGridSameRatioCovA\ and
\simGridSameRatioCovB, a gap of \simGridSameRatioGapSE\ combined Monte Carlo
standard errors; and holding the register at \simGridFixedK\ levels while
moving the training half from \simGridFixedNa\ rows to \simGridFixedNb\ moves
coverage from \simGridFixedCovA\ to \simGridFixedCovB, so $n$ carries signal
of its own. The calibration of Section~\ref{sec:regions} is a function of the ratio
alone and therefore inherits this residual, which
Section~\ref{sec:lim} states.

\textbf{Where this corpus sits, and what is done about it.} The median
log--target pair here has $K/n = \corpusRatioMedian$, the range runs
\corpusRatioMin\ to \corpusRatioMax, and \nPairsRatioAboveTenth\ of \nPairs\
pairs are above a tenth. Most of this corpus therefore sits at or beyond the
ratio at which coverage starts to fall, and a nominal band on those pairs is
\emph{narrower} than a $95\%$ interval should be --- the anti-conservative
direction, which inflates the count of resolved cells. The plane is dense
enough to calibrate against, and Section~\ref{sec:regions} does.

\textbf{Block length.} \nSimBlockLengths\ lengths are examined --- half the
rule $\lceil n^{1/3} \rceil$, the rule, and twice it. Coverage moves with the
length, spanning \coverageBlockSpreadMax\ within a world, of the order of two
standard errors; but the comparison the paper rests on is unchanged at every
length, the basic interval beating the percentile interval on the sparse
register by at least \blockSparseGapMin\ at all \nSimBlockLengths.

\paragraph{What the simulation does not establish}
It certifies the constructions on \nWorlds\ worlds, a plane of
\nSimGridCells\ $(n, K)$ cells and \nSimBlockLengths\ block lengths, and on
nothing else. The stationarity and mixing the moving-block bootstrap needs are
assumed and not tested; the split point is fixed within a draw, so
split-position uncertainty is outside every interval reported here; and the
drifting world shows what happens when the estimator is not consistent for the
quantity a reader cares about.
